\subsection{Varieties of groups}

5.12.1. Let G be a group. A manifold structure on G is said to be compatible with the group structure of G if the map \((x, y) \mapsto xy\) of \(G \times G\) into G is a morphism. The map \(x \mapsto x^{-1}\) is then a morphism from G to itself. The set G, equipped with its group structure and its manifold structure, is called a group manifold (“of class \(C^{r}\)” if one wishes to be precise), or again a Lie group. If \(r = \omega\) we also say analytic group. If K = R (resp. C, \(Q_{p}\) ), one also says real (resp. complex, p-adic) Lie group. Every group variety is pure. A homomorphism of group varieties (or simply homomorphism) is any map from one group variety to another that is both a group homomorphism and a morphism of varieties.

If G is a group variety, the topological structure underlying the variety structure of G makes it a metrizable and complete topological group; it is locally compact if K is locally compact and G is of finite dimension.

5.12.2. Examples:

\begin{enumerate}
\def\labelenumi{(\roman{enumi})}
\item
  If V is a Banach space, the canonical manifold structure of V is compatible with its commutative group structure.
\item
  Let A be a complete normed K-algebra with a unit element, and let A* be the group of invertible elements of A. It is an open subspace of A, and the manifold structure induced on this open set by the canonical manifold structure of the K-vector space A is compatible with the group structure on A*. In particular, taking A to be the algebra \(\mathcal{L}(E)\) of continuous endomorphisms of a Banach space E; the group \(A^{*}\) consists of the automorphisms of the topological vector space E; we denote by \(\mathbf{GL}(E)\) the group variety thus defined. When \(A = M_{n}(K)\) , we obtain a group variety structure on \(\mathbf{GL}(n, K)\) .
\item
  If \(G_{1},\ldots,G_{n}\) are group varieties, the product group \(G=G_{1}\times\cdots\times G_{n}\) is a group variety, when endowed with the product variety structure of those of \(G_{i}\) .
\end{enumerate}

5.12.3. Let G be a group variety, let H be a topological group, and let \(f:H\to G\) be a continuous homomorphism satisfying condition (QR) of No. 5.8.1. The inverse image manifold structure of that of G by \(f\) then makes H a group variety. This applies in particular when H is a subgroup of G which is a subvariety (resp. quasi-subvariety); such a subgroup is called a group subvariety (resp. group quasi-subvariety) of G; it is necessarily closed in G.

If \(H_{i}(i=1,\ldots,n)\) is a group subvariety of \(G_{i}\) , then \(H_{1}\times\cdots\times H_{n}\) is a group subvariety of \(G_{1}\times\cdots\times G_{n}\) .

5.12.4. Let G be a group variety and let H be a group subvariety of G. The equivalence relation \(x^{-1}y \in H\) is regular, which allows one to endow the space G/H of left classes xH with a variety structure called the quotient of that of G. The canonical map \((g, x) \mapsto g.x\) of \(G \times (G/H)\) into G/H is a morphism. Analogous results hold for the homogeneous space \(H\backslash G\) of right classes Hx. If H is normal, the variety structure of G/H is compatible with its group structure.

5.12.5. Let G be a group variety and let X be a variety. A left action of the group G on the variety X is a morphism \((s, x) \mapsto sx\) of \(G \times X\) into X such that

\[
\begin{array}{l l} s (t x) = (s t) x & \text { if } \quad s, t \in G, x \in X \\ e x = x & \text { if } \quad x \in X (e \text { being   the   identity   element   of } G). \end{array}
\]

One also says that the group variety G operates on the left on X; right action laws are defined analogously.

Let \(x \in X\) and let \(G_x\) be its stabilizer in \(G\) . Suppose that the map \(g \mapsto g \cdot x\) is a subimmersion (a hypothesis automatically satisfied if the characteristic of \(K\) is 0 and \(X\) is of finite dimension). Then \(G_x\) is a group subvariety of \(G\) and the map from \(G/G_x\) to \(X\) obtained by passing to the quotient from \(g \mapsto g \cdot x\) is an immersion. If moreover the orbit \(G \cdot x\) of \(x\) is locally closed and if the topology of \(G\) admits a countable base, \(G \cdot x\) is a subvariety of \(X\) and the map \(G/G_x \to G \cdot x\) is an isomorphism of varieties.

